\section{The seven-eighths bound}\label{sec:conclusion}

We complete the contradiction assumed in Part~II.  Recall from
Equations~\eqref{eq:part-II-bootstrap} and~\eqref{eq:part-II-bin-ceiling}
that
\[
 0<\Delta=\beta_*-\frac78\le\frac1{24},\qquad
 \kappa=\frac34+2\Delta\le\frac56,\qquad
 \delta=2a-1\le\kappa
\]
for every retained bin.  The compensated low estimate is
Proposition~\ref{prop:probe-low}. For the same physical probe, we first
normalize the principal term of its high expansion, then bound the remaining
rows and verify the target-independent margins required by
Proposition~\ref{lem:continuation-criterion}.

The geometry and Mellin exponent are
\[
 h=\frac{13}{16},\qquad \ell=\frac16,\qquad
 l_x=\frac{17}{48},\qquad l_y=\frac{23}{48},\qquad
 h=1-l_x+\ell,
\]
\[
 C(s)=C_{\mathrm{II}}(s)=s-\frac{11}{16},
 \qquad C(7/8)=\frac3{16}.
\]
These are the values in Equations~\eqref{old-eq:5.1}
and~\eqref{eq:part-II-mellin-target}.  Write \(s\) for the local variable
called \(x\) in Section~\ref{sec:local-compensation}; the amplitude ratio
\(x=q/\delta\) below is a different real number.

Use Definition~\ref{def:stage-high-data} with
\[
 \mathscr I_\eta(Z)=I_{\eta,\mathrm{modified}}(Z),\qquad
 \mathfrak H_{\eta,u,Z}(s,w,z)
     \ \text{as in Equation~\eqref{eq:holomorphic-selected-tuple}}.
\]
In the shared analytic lemmas take \(\sigma_0=7/8\).
Section~\ref{sec:local-compensation} verifies every part of these data:
Equation~\eqref{eq:compensated-poisson-identity} is the absolutely
convergent high identity for the independently defined finite expression
in Equation~\eqref{eq:compensated-probe-definition}, and the full
correction is holomorphic in the two required Euler regions and satisfies
Equation~\eqref{eq:stage-all-height-majorant}.  The excluded set, physical
row masks, and calibration are unchanged.  In particular there is no
additional Euler factor outside this full correction.  Its scalar
denominator is \(L_F^S(s,\eta\overline{\chi_\bullet(u)})\), and its
Gaussian is \(\Phi(s+z-1)\).

In the estimates below, \(0<e<10^{-3}\) is an admissible detector width.
Its final choice, together with the remaining real losses, is made in
the concluding order of choices.

\subsection{The principal normalizer}\label{sec:principal-signal}

We verify the two correction hypotheses of
Lemma~\ref{lem:principal-residue-interface}.  First,
Equation~\eqref{eq:principal-local-tuple-bound} gives
\[
 |\mathfrak H_{\eta,1,Z}(s,w,z)|\ll Z^{\ell\Re z}
\]
uniformly in all imaginary parts on every fixed real box in the second
Euler region.  In particular this holds on the entire rectangle in
Equation~\eqref{eq:principal-correction-bound}, with height degree zero.
This is a bound for the full tuple correction.

For the residue value, define
\[
 S_i(Z)=\sum_{p\in\mathcal P_i(Z)}W_i(q_p/P_i)q_p^{-5/6},
 \qquad P_i=Z^{\ell_i}.
\]
Lemma~\ref{lem:ray-prime-normalizer}, including its assertion about deletion
of a fixed finite set, gives
\[
 S_i(Z)\sim\frac{P_i^{1/6}}{|T|\log P_i}
       \int_0^\infty W_i(y)y^{-5/6}\,dy.
\]
Each leading constant is positive.  Since the slot count is fixed,
all \(S_i(Z)\) are positive for every sufficiently large \(Z\), with
a common lower threshold allowed to depend on the fixed data.  Set
\begin{equation}\label{eq:principal-scalar}
 A_T(Z)=(-1)^K Z^{-\ell/6}\prod_{i=1}^K S_i(Z).
\end{equation}
Using \(\log P_i=\ell_i\log Z\) and \(\sum_i\ell_i=\ell\), we obtain
\[
 A_T(Z)\sim
 \frac{(-1)^K}{|T|^K(\log Z)^K}
 \prod_{i=1}^K\left\{\frac1{\ell_i}
       \int_0^\infty W_i(y)y^{-5/6}\,dy\right\}.
\]
Consequently \(A_T(Z)\ne0\) on that range and
\[
 |A_T(Z)|\asymp_{T,K,(\ell_i,W_i)}(\log Z)^{-K},
 \qquad |A_T(Z)|^{-1}\ll_\epsilon Z^\epsilon
\]
for every \(\epsilon>0\).  The lower threshold may depend on the target's
finite excluded set; no uniform prime asymptotic in a moving ray conductor
is being used.

On the principal second Euler region, Section~\ref{sec:local-compensation}
defines \(\mathcal B_p=G_p/H_p\) with \(H_p\ne0\), and
Equation~\eqref{eq:principal-local-approximation} gives
\(\mathcal B_p=-1+O(q_p^{-7/8})\), uniformly in all imaginary parts and
target unit phases.  Apply the principal factorization in
Equation~\eqref{eq:principal-slot-factorization} at \(w=1,z=1/6\), where
it remains valid even if the original quotient by \(P_p\) was undefined.
It gives
\[
 \mathfrak H_{\eta,1,Z}(s,1,1/6)
 =H_\eta(s)\prod_{i=1}^K
       \left(\sum_{p\in\mathcal P_i(Z)}
           W_i(q_p/P_i)q_p^{-5/6}\mathcal B_p\right),
\]
where \(H_\eta\) is the function in Equation~\eqref{eq:shared-principal-data}
for this same excluded set \(S\).  By nonnegativity of \(W_i\), its \(i\)th
slot equals
\[
 -S_i(Z)\{1+\rho_i(s)\},\qquad
 |\rho_i(s)|\ll P_i^{-7/8}
 \quad(\Re s=\beta_*+e,\ \Im s\in\mathbb R).
\]
Indeed the sum of the absolute local errors is at most a fixed multiple of
\(P_i^{-7/8}S_i(Z)\).  Choose any pretarget number
\[
 0<\kappa_P<\frac78\min_i\ell_i.
\]
Since \(K\) is fixed, \(\mathcal R_{\eta,Z}(s):=\prod_i(1+\rho_i(s))-1\)
satisfies \(|\mathcal R_{\eta,Z}(s)|\ll Z^{-\kappa_P}\) on that entire line.
Thus the exact residue correction is
\[
 \mathfrak H_{\eta,1,Z}(s,1,1/6)
 =H_\eta(s)Z^{\ell/6}A_T(Z)\{1+\mathcal R_{\eta,Z}(s)\},
\]
which is Equation~\eqref{eq:principal-correction-residue} with
\(A_\eta=A_T\) and \(\mu=\kappa_P\).  This factorization is asserted
only at the principal residue; it is not the correction used for general
contour moves.

Let \(c_S>0\) be the scalar in Equation~\eqref{eq:shared-principal-data},
whose positivity is part of Lemma~\ref{lem:principal-residue-interface}.
For sufficiently large \(Z\), define the normalized physical probe
\[
 J_{\mathrm{II},\eta}(Z)
   =\frac{I_{\eta,\mathrm{modified}}(Z)}{c_S A_T(Z)}.
\]
This is distinct from \(J_{\mathrm I,\eta}\), and both its numerator and
normalizer are independent of the analysis height \(T_1\).
Proposition~\ref{prop:probe-low} and the subpower bound for \(A_T^{-1}\)
give, for every \(\epsilon>0\),
\begin{equation}\label{eq:part-II-normalized-low}
 |J_{\mathrm{II},\eta}(Z)|
 \ll_{\eta,\epsilon} Z^{C(7/8)+\epsilon}.
\end{equation}
The same scalar is used for the high comparison.

Choose the fixed cutoff as required for
Equation~\eqref{eq:shared-principal-nonvanishing}.  The associated
\(H_\eta\) is holomorphic on \(\Re s>7/8\), satisfies the contraction
there, and uses this same \(S\).  Let \(\mathscr P_{\mathrm{II},\eta}\)
denote the \(u=1\) term of Equation~\eqref{eq:compensated-poisson-identity},
and set
\[
 f_{\mathrm{II},\eta}(Z)=\frac1{2\pi i}\int_{(2)}
 Z^{C(s)}e^{(s-5/6)^2}\frac{H_\eta(s)}{L_F^S(s,\eta)}\,ds.
\]
All hypotheses of Lemma~\ref{lem:principal-residue-interface} have now
been verified with \(\sigma_0=7/8\).  Its remainder estimate is
\[
 \begin{split}
 \left|\frac{\mathscr P_{\mathrm{II},\eta}(Z)}{c_S A_T(Z)}
       -f_{\mathrm{II},\eta}(Z)\right|
 \ll{}&
 Z^{C(\beta_*)+(1+h)e-m_w+\epsilon}
 +Z^{C(\beta_*)+e-m_z+\epsilon}\\
 &+Z^{C(\beta_*)+e-\kappa_P+\epsilon},
 \end{split}
\]
where the two geometric margins are
\begin{equation}\label{old-eq:5.14}
 m_w:=\frac{l_y}{20}=\frac{23}{960},\qquad
 m_z:=\frac h{600}=\frac{13}{9600}.
\end{equation}
Lemma~\ref{lem:principal-residue-interface} includes the residue paths,
the scalar Jacobian, and
the rightward shift of the main integral.  In particular the last term
is the error estimated on its original global line; it is not part of
\(f_{\mathrm{II},\eta}\).  Both this signal and the physical probe are
independent of \(T_1\).

\subsection{The compensated high exponent}

The low estimate and the principal-row comparison now concern the functions
$J_{\mathrm{II},\eta}$ and $f_{\mathrm{II},\eta}$ required by the continuation
criterion. It remains to prove, for a common $\sigma>0$,
\[
 \left|\frac{I_{\eta,\mathrm{modified}}(Z)
                -\mathscr P_{\mathrm{II},\eta}(Z)}{c_S A_T(Z)}\right|
 \ll_\eta Z^{C(\beta_*)-\sigma}.
\]
We estimate these nonprincipal contributions first relative to the raw
scale $Z^{C(7/8)}$. Division by $c_S A_T(Z)$ costs an arbitrarily small
power, reserved once in the final choice of margins.


\begin{lemma}[A high exponent for one row bin]\label{lem:high-bin}
Let \(0<d_{\min}<d_{\max}<\infty\), let \(U=Z^d\) with
\(d_{\min}\le d\le d_{\max}\), and let
\(\mathcal B\) be the finite set of retained nonprincipal physical rows
\(q_u\asymp U\) in one fixed dynamic bin \((i,a)\).  Put
\(\delta=2a-1\).  Assume the bin and height hypotheses of
Lemmas~\ref{lem:buffered-bins}, \ref{lem:detector-dyads}, and
\ref{lem:bin-prime-bound}, with
\(0<e<10^{-3}\), \(T_1=Z^\tau>2\), \(0<\tau\le d_{\min}/100\), and the cumulative
allocation in Equation~\eqref{eq:stage-height-allocation}.
The bounded real ranges, the positive losses \(e,\vartheta,\epsilon\), and the
slot system are fixed before the target.

For each retained tuple of external parameters and each subset of error
slots, partition \(\mathcal B\) pointwise into the amplitude sets of
Section~\ref{sec:amplitudes}, and into the witness subdivisions of
Section~\ref{sec:selection} when a witness count is used.  Suppose each
such set \(\mathcal C\) has main-slot mean \(q\in[0,\delta/2]\), as in
Equation~\eqref{old-eq:6.6}, and
\[
 \#\mathcal C\ll_{\mathcal A,\epsilon}
 U^{R+\epsilon}(1+T_1)^{A_{\mathcal A}},
\]
uniformly in the retained tuple and moving labels.  The occurring \(R,q\)
range over a fixed bounded set and may depend on the pointwise set, \(d,a\),
and \(\beta_*\), but not otherwise on the target.  For each occurring pair
define
\begin{equation}\label{eq:common-high-exponent}
 \begin{split}
 E(d)
 &=a-\frac78+h\left(\frac{17}{50}-\frac16\right)
       -a l_y-(1-a)\ell-(\delta/2-q)\ell
       +d\left(R+\frac\delta2-\frac{17}{50}\right)\\
 &=C_0+\frac23\delta+\frac q6-h(1-R)
       +(d-h)\left(R+\frac\delta2-\frac{17}{50}\right),
       \qquad C_0=-\frac1{48}.
 \end{split}
\end{equation}
If \(\mathcal B\ne\varnothing\), let \(E_{\max,Z}(d)\) be the supremum
of these values over all pointwise sets at all retained tuples.

On the central contours
\[
 \Re s=a+16e,\qquad \Re w=1-a-6e,\qquad \Re z=\frac{17}{50},
\]
the contribution of the original fixed dynamic-bin sum is bounded by
\[
 Z^{C(7/8)+E_{\max,Z}(d)+O(e+\vartheta+\epsilon)}
 (1+T_1)^{A'_{\mathcal A}},
\]
where \(\vartheta\) is the amplitude-bin width and
\(A'_{\mathcal A}<\infty\) is uniform in moving labels.  An empty bin
contributes zero.  The constant in the \(O(e+\vartheta+\epsilon)\) term
depends only on the bounded real ranges and the fixed slot system, not on
the target.  No separate integral of an individual pointwise set is
asserted.  After fixing \(\tau>0\), all discarded external pieces and joins
are \(O_{\mathcal A,N}(Z^B T_1^{-N})\), for a fixed \(B\) and every fixed
\(N\).
\end{lemma}

\begin{proof}
We verify the central hypothesis of
Lemma~\ref{lem:bin-contour-accounting}, which supplies the contour
move and its tails for the full correction just specified.  Its bin
ceiling holds by Equation~\eqref{eq:part-II-bin-ceiling}.  The physical
\(\Im z\) coordinate is included in its height allocation because it
enters the prime profiles.  The other witness, dyadic, and prime-annulus
coordinates are the fixed finite list used in the estimates of
Section~\ref{sec:selection}.  Their complete pointwise bounds include the
discarded auxiliary integrals on global or absolute lines.  Thus the
single cumulative allocation in Equation~\eqref{eq:stage-height-allocation}
applies, without renewing an allowance at a later estimate.

Only on the retained contours, Proposition~\ref{prop:probe-errors} gives
the decomposition in Equation~\eqref{eq:dynamic-compensated-decomposition}.
For \(I\subseteq\{1,\ldots,K\}\), its summand is
\[
 \mathfrak H_{\eta,u,Z}^{(I)}
 =\mathcal H_{\eta,u}(s,w,z)
       \prod_{i\in I}\mathcal D_i(u)
       \prod_{i\notin I}\mathcal Q_i(u;z).
\]
The required reflected primitive numerator bound follows from the
buffered estimate for the numerator presentation and its conjugate,
together with Lemma~\ref{lem:deleted-euler-factors}, as verified before
Proposition~\ref{prop:probe-errors}.  The retained \(w\) heights lie
in its allowed set.  Every retained numerator is nonprincipal.

Fix one amplitude set at one retained tuple.  For a main slot,
\(\mathcal Q_i=-P_i^{z-1/2}Q_i\) and
\(\lvert Q_i\rvert\le P_i^{g_i+\vartheta}\).  Assign \(g_i=0\) to
an error slot.  Equation~\eqref{old-eq:5.13e} bounds all error slots
in \(I\) jointly with the numerator; its proof uses the actual
conductor deficit in Equation~\eqref{old-eq:5.13d} simultaneously
for the distinct strict ramified labels.  It is not a separate
numerator allowance for each slot.  Since
\(\mathcal H_{\eta,u}\ll_\epsilon U^\epsilon\), multiplication of
these estimates gives, after choosing the preliminary powers,
\[
 \begin{split}
 \left|L^S(w,\chi_\bullet(u))
       \mathfrak H_{\eta,u,Z}^{(I)}(s,w,z)\right|
 \ll{}&U^{\delta/2+O(e)+\epsilon}(1+T_1)^{A_1}\\
 &\cdot
 Z^{\ell(17/50-1/2)+q\ell+O(e+\vartheta+\epsilon)}.
 \end{split}
\]
Here \(A_1<\infty\) is a fixed height order for the target and slot
system, and \(\sum_i\ell_i g_i=q\ell\), including the assigned zeros
for the errors.  No lower bound on an error slot has been used.

Multiplying by the assumed cardinality proves
Equation~\eqref{eq:stage-central-bound} with \(g=q\ell\), after taking
\(\varepsilon_c\) to be a fixed sufficiently large multiple of
\(e+\vartheta+\epsilon\).  This multiple depends only on the fixed slot
system and the bounded real ranges.  There are \(2^K\) error subsets.
The number of amplitude vectors is fixed, and the witness dyadic choices
contribute only a fixed power of \(1+\log Z\).  Thus the required
pointwise multiplicity bound holds.  These sets are formed only after
Lemma~\ref{lem:bin-contour-accounting} has moved the original dynamic-bin
sum using
\(\mathfrak H\).  They are not used for continuation and are not
integrated separately.

Apply Lemma~\ref{lem:bin-contour-accounting} with \(\sigma_0=7/8\)
and \(g=q\ell\).  Its exponent in
Equation~\eqref{eq:stage-high-exponent} is
\[
 a-\frac78+h\left(\frac{17}{50}-\frac16\right)-a l_y-\frac{\ell}{2}
       +q\ell+d\left(R+\frac\delta2-\frac{17}{50}\right).
\]
Because
\(-(1-a)\ell-(\delta/2-q)\ell=-\ell/2+q\ell\),
this is the first line of Equation~\eqref{eq:common-high-exponent}.
Substituting the geometry and \(a=(1+\delta)/2\) gives its second line.
In particular \(q\ell=q/6\) is a base-\(Z\) exponent, not \(dq/6\).
The explicit errors in Equation~\eqref{eq:stage-central-contribution}
are \(O(e+\vartheta+\epsilon)\) on the bounded \(d\)-range after the
remaining preliminary powers are chosen.  Its supremum is exactly the
one in the statement.  Lemma~\ref{lem:bin-contour-accounting} also gives
the stated external tails
with a scale degree fixed before \(N\).
\end{proof}

The floor bin \(a=51/100\) has \(\delta_0=1/50\) and requires no
zero witness.  The ideal count \(R=1\), together with
\(q\le\delta_0/2\), gives
\begin{equation}\label{eq:floor-bound}
 E(h)\le C_0+\frac34\delta_0=-\frac7{1200}<0.
\end{equation}
Its frequency slope \(R+\delta_0/2-17/50\) is positive, so this
bounds every \(d\le h\).  A small extension above \(h\) will be
controlled below.

\subsection{Small and large row norms}\label{sec:tails}

Set \(d_{\min}=1/100\).  Lemma~\ref{lem:compensated-absolute-local-bounds}
bounds the positive sum of the full selected tuples, with every \(G_p\)
retained before taking absolute values.  On the small-row lines
\[
 \Re s=\beta_*+e,\qquad \Re w=1/2,\qquad \Re z=17/50
\]
Equation~\eqref{eq:absolute-local-tuple-bound} gives
\[
 |\mathfrak H_{\eta,u,Z}(s,w,z)|
 \ll_\epsilon U^\epsilon\prod_iP_i^{17/50}
 =U^\epsilon Z^{17\ell/50},
\]
uniformly in all imaginary parts.  This is
Equation~\eqref{eq:small-correction-bound}, with height degree zero.
It remains valid at zeros of individual local factors.  The data for
Lemma~\ref{lem:outer-row-tails} were already verified above, so it
applies to every physical row \(u\ne1\) in these dyads.  Its
classification includes nontrivial unit numerators and permits a principal
denominator; no detector witness or selected moment is used here.

The \(d\)-independent part of the relative exponent is
\begin{equation}\label{eq:small-frequency-margin}
 h\left(\frac{17}{50}-\frac16\right)-\frac{l_y}{2}
 =-\frac{79}{800}.
\end{equation}
The row exponent in Equation~\eqref{eq:small-row-bound} is
\(63/50<2\).  Thus the small dyadic sum in
Equation~\eqref{eq:small-row-sum}, measured relative to \(C(\beta_*)\),
has exponent at most
\[
 -\frac{79}{800}+2d_{\min}+O(e+\epsilon)
 =-\frac{63}{800}+O(e+\epsilon).
\]
It is negative after the adjustable real losses are made sufficiently
small.  Write \(m_{\rm small}:=63/800\) for this fixed error-free saving.

For the large rows, fix \(z_\infty>2\).  On
\((\Re s,\Re w,\Re z)=(2,2,z_\infty)\), the same
Equation~\eqref{eq:absolute-local-tuple-bound} gives
\[
 |\mathfrak H_{\eta,u,Z}(s,w,z)|
 \ll_\epsilon U^\epsilon Z^{\ell z_\infty}
\]
for every physical row and all imaginary parts.  This verifies
Equation~\eqref{eq:large-correction-bound}.  Lemma~\ref{lem:outer-row-tails}
therefore gives, with \(B_0=l_x/2+1+l_y=53/32\),
\[
 |\mathscr R_\eta(U;Z)|
 \ll Z^{B_0+hz_\infty+\epsilon}U^{1+\epsilon-z_\infty},
\]
and, for every fixed \(\zeta>0\) after choosing
\(0<\epsilon<z_\infty-1\),
\[
 \sum_{U>Z^{h+\zeta}}|\mathscr R_\eta(U;Z)|
 \ll
 Z^{B_0+(h+\zeta)(1+\epsilon)-\zeta z_\infty+\epsilon}.
\]
Here \(\mathscr R_\eta\) is the row contribution of
Lemma~\ref{lem:outer-row-tails} for the
present physical expression.  A fixed sufficiently large \(z_\infty\)
makes the last exponent as negative as required.  It is chosen after
\(\zeta\) but before the target.  These invocations use the quotient-free
tuple bounds in Section~\ref{sec:local-compensation}, not the central
main/error factorization.

\subsection{The endpoint inequality}\label{sec:adaptive-endpoint}

We next bound \(E(h)\), first without adjustable losses or height
factors.  They will be restored with a quantified order of choices.
Recall
\[
 \kappa=\frac34+2\Delta,\qquad \alpha=\frac56.
\]
At the live upper endpoint \(\delta=\alpha\), take \(t=3/2\) in
Proposition~\ref{prop:detector-witness}.  Then
\(r\ge1-O(\epsilon)\).  Apply the no-slot amplified estimate
in Equation~\eqref{eq:no-slot-inverse-count} on both sides of
\(r=1\).  For \(r\ge1\), its exponent is
\[
 1-\alpha+(\alpha-\delta)r=1-\delta.
\]
For \(r<1\) in the stated \(O(\epsilon)\) neighborhood, the same
estimate uses \(e(r)=1\), so its exponent is
\(1-\delta r\le1-\delta+O(\epsilon)\).  The fixed amplification
margin is independent of this neighborhood.
Thus no selected primes are needed and the ideal row exponent is
\(R=1-\delta\).  Using \(q\le\delta/2\) in
Equation~\eqref{eq:common-high-exponent} gives
\begin{equation}
 E(h)\le C_0+\left(\frac34-h\right)\delta+O(\epsilon)
       =-\frac1{48}-\frac{\delta}{16}+O(\epsilon)<0.
 \label{eq:large-delta-endpoint}
\end{equation}

Now assume \(1/50<\delta<\alpha\).  Put \(x=q/\delta\in[0,1/2]\),
and use \(D_x,P_x,R_{\rm short},L\) from
Proposition~\ref{prop:detector-counts}.  The parameter \(\kappa\)
remains the dynamic value \(3/4+2\Delta\), not the upper bound
\(5/6\).  Thus the actual plain capacity used in that Proposition is
\[
 z_P(m)=\frac{1-2m}{6\kappa}
       =\frac{1-2m}{9/2+12\Delta},
\]
as in Equation~\eqref{eq:plain-zero-loss-capacity}.  The comparison in
Equation~\eqref{eq:capacity-comparison} is for this exact capacity and
raises the row exponent by at most \(\Delta/4\), apart from requested
small losses.  Define
\begin{equation}
 \mathcal J=(\alpha-\delta)D_x+\delta P_x,\qquad
 t=1+\frac{\delta P_x}{2\mathcal J},\qquad
 R_*=L(t).
 \label{eq:target-cutoff}
\end{equation}
The denominator is positive.  Indeed
\[
 \frac{37}{18}\le D_x\le3,\qquad
 \frac79\le P_x\le2,\qquad
 D_x-P_x=1+x-\frac{8x^2}{9}>0,
\]
and consequently
\[
 \frac{35}{54}\le\mathcal J\le\frac52.
\]
Because \(0<\delta<\alpha\), one has
\(\mathcal J>\delta P_x>0\), so \(1<t<3/2\).
Furthermore,
\[
 D_x\{R_{\rm short}(t)-L(t)\}
 =\delta P_x(3/2-t)-(\alpha-\delta)D_x(t-1)
 =\frac{\delta P_x}{2}-\mathcal J(t-1)=0.
\]
Thus \(R_*=R_{\rm short}(t)=L(t)\) and this \(t\) balances the two counts.
The formulas also have the stated closed-endpoint values: at
\(\delta=\alpha\), one has \(\mathcal J=\alpha P_x\), \(t=3/2\),
and \(R_*=R_{\rm short}(3/2)=L(3/2)=1-\delta\).  The corresponding crossing
has \(r_*(3/2)=1\), so its inverse capacity is zero.  The preceding
equality argument uses Equation~\eqref{eq:no-slot-inverse-count} at
that boundary, not a marked estimate with a vanishing margin.

\begin{lemma}[Compensated endpoint certificate]\label{lem:balanced-endpoint}
For \(0\le\delta\le5/6\) and \(0\le x\le1/2\), define
\(\mathcal J,t,R_*\) by Equation~\eqref{eq:target-cutoff} where
\(\mathcal J>0\), and let \(E_*\) be \(E(h)\) with
\(q=x\delta\) and \(R=R_*\).  Then
\[
 -E_*\ge\frac{49}{440640}>\frac1{10000}.
\]
In particular the bound is uniform on the range
\(1/50<\delta\le5/6\), including the closed endpoint.
\end{lemma}

\begin{proof}
Set \(y=1/2-x\in[0,1/2]\) and \(v=51+41y\).  To verify the
calculation, write
\[
 p_y=7+18y+8y^2,\qquad
 j_y=185+170y+(-138+12y+96y^2)\delta.
\]
Then
\[
 D_x=\frac{37+34y}{18},\qquad
 P_x=\frac{p_y}{9},\qquad
 \mathcal J=\frac{j_y}{108}.
\]
From \(R_*=1-\delta+(\alpha-\delta)\delta P_x/(2\mathcal J)\)
and Equation~\eqref{eq:common-high-exponent},
\[
 -E_*=
 \frac{1+(3+8y)\delta}{48}
 -\frac{13}{32}\,
   \frac{(5/6-\delta)\delta P_x}{\mathcal J}.
\]
Multiplying out gives the explicit quadratic
\[
 \begin{split}
 10368\mathcal J(-E_*)
 &=2j_y[1+(3+8y)\delta]
     -468(5/6-\delta)\delta p_y\\
 &=10(37+34y)
   -8(237+377y+26y^2)\delta\\
 &\hspace{12mm}
   +48(51+131y+94y^2+32y^3)\delta^2.
 \end{split}
\]
Multiplication by \(v\) and completion of the square yields
\begin{equation}
 \begin{aligned}
 10368v\mathcal J(-E_*)
 ={}&(3+5y)\bigl\{(4v\delta-79)^2+49\bigr\}\\
 &+4y\bigl\{
 4v\delta[(1+3y)(15+32y)\delta+9-13y]
       +265+3485y\bigr\}.
 \end{aligned}
 \label{eq:balanced-endpoint-certificate}
\end{equation}
This identity can also be checked by coefficients: expansion of
its right side as a polynomial in \(\delta\) gives respectively
\[
 \begin{split}
 &10v(37+34y),\\
 &-8v(237+377y+26y^2),\\
 &48v(51+131y+94y^2+32y^3),
 \end{split}
\]
which are the constant, linear, and quadratic coefficients in
the preceding display multiplied by \(v\).

Every term on the right of
Equation~\eqref{eq:balanced-endpoint-certificate} is nonnegative:
\(y,\delta\ge0\) and \(9-13y\ge5/2\).
The first line is at least \(49(3+5y)\), while
\(v=51+41y\le17(3+5y)\).  Hence
\[
 -E_*\ge\frac{49}{176256\mathcal J}
       \ge\frac{49}{440640}>\frac1{10000},
\]
using \(\mathcal J\le5/2\).
\end{proof}

The lower bounds for \(D_x\) and \(\mathcal J\) bound all derivatives
of the crossing and cutoff on the compact parameter ranges.
The saturation losses are uniform because \(\delta\ge1/50\),
and the inverse selection has the explicit supply and width margins
proved in Proposition~\ref{prop:detector-counts}.  At a zero
capacity its unweighted bound differs from \(1-\delta\) by at
most a constant times the chosen capacity decrement.  Since
\(R_*\ge1-\delta\), these cases obey the same comparison.
Thus all detector, bin, and rounding errors have a total
\(O(\epsilon)\) with a target-independent constant once their
individual requested losses are at most \(\epsilon\).
Let \(E_{\rm actual}(h)\) denote the scale exponent relative to
\(C(7/8)\) after the observed row bound
\(R\le R_*+\Delta/4+\epsilon_{\rm row}\) is inserted.
For now the explicit height factors remain outside this exponent.
The coefficient of \(R\) in Equation~\eqref{eq:common-high-exponent}
is \(d\), so at \(d=h\) the capacity replacement contributes
\(h\Delta/4\).  Denote by \(\epsilon_{\rm total}\) the sum of
\(h\epsilon_{\rm row}\) and all other adjustable real losses; the
preceding uniformity makes it arbitrarily small with a
target-independent coefficient.  The certificate gives the
comparison actually needed:
\begin{equation}
 \begin{split}
 E_{\rm actual}(h)-\Delta
 &\le-\frac{49}{440640}-\left(1-\frac h4\right)\Delta
       +\epsilon_{\rm total}\\
 &=-\frac{49}{440640}-\frac{51}{64}\Delta
       +\epsilon_{\rm total}.
 \end{split}
 \label{eq:adaptive-endpoint}
\end{equation}
The subtraction of \(\Delta=C(\beta_*)-C(7/8)\) is essential:
the certificate is not a claim that \(E_{\rm actual}(h)<0\) for
every \(\Delta\).  When the height factors are converted to a
reserved exponent allowance below, that allowance is added to
\(\epsilon_{\rm total}\).

\subsection{Frequency ranges}

For \(1/2\le d\le h\), the available length is
\[
 \frac{\ell}{d}\ge\frac8{39}>\frac15>\frac7{37}.
\]
Passing from base \(Z\) to base \(U\) changes a slot length from
\(\ell_i\) to \(\ell_i/d\le2\ell_i\), so one mesh chosen with
this factor of two satisfies the moment requirement throughout the
interval.  If
\[
 0<\zeta<5\ell-h=\frac1{48},
\]
then for \(h<d\le h+\zeta\) the supply remains greater than
\(1/5\), whose margin above \(7/37\) is \(2/185\).

For \(1/50<\delta<\alpha\), use the ideal row upper exponent
\(R=R_*+\Delta/4\).  For \(\delta=\alpha\), use
\(R=1-\delta\).  Both are at least \(1-\delta\), so
\[
 R+\delta/2-\frac{17}{50}
 \ge\frac{33}{50}-\frac{\delta}{2}\ge\frac4{25}>0.
\]
The floor has a positive slope as well.  Thus \(d=h\) bounds
all these exponents for \(d\le h\).  The cost of extending to
\(h+\zeta\) is at most \(2\zeta\): indeed
\[
 R_*\le1-\delta+\frac{\alpha-\delta}{2}\le\frac{17}{12},
 \qquad
 R_*+\Delta/4\le\frac{139}{96},
\]
and \(139/96+1/2-17/50<2\).  The slopes at \(\delta=\alpha\)
and floor cases are smaller than two.

For \(d_{\min}\le d\le1/2\) and \(\delta\le\alpha\), choose \(t=1\)
and use no selected primes.  Proposition~\ref{prop:detector-counts}
gives the ideal exponent \(1-2\delta/3\) outside the floor.
Use the slightly larger common exponent
\[
 R=\frac{76}{75}-\frac23\delta.
\]
At \(\delta=1/50\) it equals one, so it also covers the floor.
Its slope is \(101/150-\delta/6>0\) on this range.
Using \(q\le\delta/2\) in
Equation~\eqref{eq:common-high-exponent} yields
\begin{equation}
 E(d)\le E(1/2)+O(\epsilon)
 \le-\frac{529}{2400}+\frac{25}{96}\delta+O(\epsilon)
 \le-\frac{49}{14400}+O(\epsilon)<0.
 \label{eq:adaptive-intermediate}
\end{equation}
The middle expression follows by substituting
\(d=1/2\), \(R=76/75-2\delta/3\), and \(q=\delta/2\);
the last inequality uses \(\delta\le5/6\).
At \(\delta=\alpha\), the count used in
Equation~\eqref{eq:large-delta-endpoint} and its positive slope
already control every \(d\le h\).
Together with Section~\ref{sec:tails}, these cases cover every
physical row norm.  Only \(d\ge1/2\) uses selected physical prime factors
in a row moment; every physical slot remains in the high correction in all
ranges.

To compare with \(C(\beta_*)\), subtract
\(C(\beta_*)-C(7/8)=\Delta\) from each ideal exponent.
In the adaptive range this is precisely the comparison in
Equation~\eqref{eq:adaptive-endpoint}; its only positive capacity
allowance is \(h\Delta/4\).  The other endpoint ranges have
nonpositive ideal exponents before subtracting \(\Delta\).
Thus all of them retain the target-independent main high margin
\[
 m_{\rm hi}:=\left(1-\frac h4\right)\Delta
             =\frac{51}{64}\Delta.
\]
The low exponent has margin \(m_{\rm lo}:=\Delta\).
The principal contour and approximation margins were established in
Section~\ref{sec:principal-signal}; we now choose all losses together.

\subsection{Order of choices and conclusion}

We finish by separating the target-independent real choices from the
target-dependent height and external test orders.  The numerical choices
remain part of this application; the final height selection is supplied by
Lemma~\ref{lem:late-height-closure}.

\begin{proposition}[Order of choices]\label{prop:parameter-order}
Under the contradiction in Equation~\eqref{eq:part-II-bootstrap}, there
are numbers
\[
 0<\omega<\Delta,\qquad \sigma>0
\]
and a fixed geometry, moment losses, capacity decrements, slot system,
amplitude width, bin width \(e>0\), and extension \(0<\zeta<1/48\),
all chosen before the target character, with the following property.
For every primitive finite-order target \(\eta\), one can then choose
its fixed arithmetic data, a positive height exponent \(\tau_\eta\),
a finite external tail order \(N_\eta\), and a lower threshold
\(Z_{0,\eta}\) such that, for \(Z\ge Z_{0,\eta}\),
\[
 |J_{\mathrm{II},\eta}(Z)|\ll_\eta Z^{C(7/8)+\omega},\qquad
 |J_{\mathrm{II},\eta}(Z)-f_{\mathrm{II},\eta}(Z)|
 \ll_\eta Z^{C(\beta_*)-\sigma}.
\]
The exponents \(\omega,\sigma\) do not depend on \(\eta\).
The cutoff \(T_1=Z^{\tau_\eta}\) occurs only in the estimates,
not in either function.
\end{proposition}

\begin{proof}
The ideal margins are
\[
 m_{\rm lo}=\Delta,\qquad m_{\rm hi}=\frac{51}{64}\Delta,\qquad
 m_{\rm high}:=\min\{m_{\rm hi},m_{\rm small}\}>0.
\]
The principal contour margins are \(m_w,m_z\) in
Equation~\eqref{old-eq:5.14}.  Reserve one half of each of
\(m_{\rm lo},m_{\rm high},m_w,m_z\).  On the compact real parameter
ranges, the coefficients multiplying all requested detector, moment,
capacity, and rounding losses are bounded independently of the target.
Indeed, \(\delta\ge1/50\), \(D_x\ge37/18\), and
\(\mathcal J\ge35/54\); the dyadic lengths are bounded, and
Proposition~\ref{prop:detector-counts} gives fixed inverse width and
supply margins.  The equality \(\delta=\alpha\) uses the separate
no-slot bound above.  Choose the moment losses and capacity decrement
so that their total cost in the high exponent is less than
\(m_{\rm high}/8\).

Let \(\eta_{\rm mesh}>0\) be the mesh supplied by
Lemma~\ref{lem:plain} for these losses and bounded real ranges.  Its
uniform assertion on \([3/4,1]\) applies in particular to the compact
closure \(\kappa\in[3/4,5/6]\) used here, and the mesh is independent
of the number of slots.  Only the eventual seminorm and height orders
may depend on a fixed slot count.  Let \(b_{\rm round}>0\) be small
enough that a squared-spike rounding loss below \(b_{\rm round}\), after
the bounded changes of exponent base, costs less than \(m_{\rm high}/8\).
Choose a fixed even integer \(K\) with
\[
 \frac{2\ell}{K}<
 \min\left\{\eta_{\rm mesh},b_{\rm round},\frac1{185}\right\},
 \qquad \ell_i=\frac{\ell}{K}\quad(1\le i\le K).
\]
These positive lengths sum to \(\ell\) and are independent of the
target and \(Z\).  Put \(P_i=P=Z^{\ell/K}\).  For \(1\le i\le K\), set
\[
 I_i=\left(1+\frac{2i-1}{2K+1},\,1+\frac{2i}{2K+1}\right)\subset(1,2)
\]
and choose a nonnegative, nonzero \(W_i\in C_c^\infty(I_i)\).
The intervals have positive gaps.  Hence their underlying prime windows
are disjoint for every \(Z\), before the ray, excluded-set, or row masks
are imposed.  Equation~\eqref{eq:compensated-probe-definition} keeps each
window at its original scale in every summand.

For each selected row range, \(U=Z^d\) has \(d\ge1/2\), so
\[
 w_i=\frac{\ell_i}{d}\le\frac{2\ell}{K}<\eta_{\rm mesh},
 \qquad
 \delta\max_iw_i\le\frac{2\ell}{K}<b_{\rm round}.
\]
The second inequality uses \(\delta\le5/6<1\).  The total available
length remains \(\ell/d\).  For \(d\le h\) it is at least \(8/39\),
exceeding \(7/37\) by \(23/1443\).  For \(h<d\le h+\zeta\) with
\(\zeta<1/48\), it exceeds \(1/5\), whose excess over \(7/37\) is
\(2/185\).  Thus every selected inverse capacity and the smaller
plain capacity are supplied.  The extra bound \(2\ell/K<1/185\)
is below half the latter gap.  A zero-capacity neighborhood, including
the equality endpoint, uses the no-slot moment.  Actual annular ratios
change nominal lengths by \(O_K(1/\log U)\), absorbed by the fixed
strict margins at a sufficiently large threshold.

The internal centered amplifier also stays separated from these windows.
Write \(\sigma_{\rm width}\) for the width decrement denoted \(\sigma\)
in the proof of Lemma~\ref{lem:plain}.  That proof uses the pool exponent
\(\ell_*=\sigma_{\rm width}/3\) and a mesh below
\(\sigma_{\rm width}/6\).  At base \(U\), every live physical slot has
norm at most \(2U^{\eta_{\rm mesh}}\), whereas the pool begins at
\(U^{\ell_*}/2\).  Their fixed exponent gap makes these sets disjoint
for all sufficiently large \(U\), uniformly in the selected \(d\)-range.
The marked inverse moment has no mesh condition.

The normalizer in Equation~\eqref{eq:principal-scalar} is nonzero for
this actual choice.  Put
\(c_i=\int_0^\infty W_i(y)y^{-5/6}\,dy>0\).  The asymptotic already
proved in Section~\ref{sec:principal-signal} specializes to
\[
 S_i(Z)\sim\frac{P^{1/6}c_i}{|T|\log P}>0,\qquad
 A_T(Z)\sim
 \left(\frac{K}{|T|\ell}\right)^K
       \left(\prod_i c_i\right)(\log Z)^{-K}>0.
\]
The sign uses even \(K\).  Any later fixed excluded set affects only
the common lower threshold, since its primes eventually leave all
these windows.  Choose
\[
 \kappa_P=\frac7{16}\frac{\ell}{K}=\frac7{96K},
 \qquad 0<\kappa_P<\frac78\min_i\ell_i=\frac7{48K},
 \qquad m_P:=\kappa_P.
\]
Thus the principal approximation has a positive margin chosen before
the target.  The number of subsets, the seminorms of the narrow windows,
and their finite height orders may depend on this fixed \(K\).

Choose the amplitude width, the remaining power losses, and \(e\) so
that their total high cost is less than \(m_{\rm high}/8\), and so that
\[
 (1+h)e+\epsilon_{\rm pr}<m_w/4,\qquad
 e+\epsilon_{\rm pr}<m_z/4,\qquad
 e+\epsilon_{\rm pr}<m_P/4.
\]
Here \(\epsilon_{\rm pr}>0\) includes the chosen small powers in the
principal remainder estimate.  Also require \(e<10^{-3}\) and the bounds
\(e_0\) in Lemma~\ref{lem:detector-dyads}.  The slot lengths have already
been fixed, so the small powers needed in the amplitude subdivision can
be chosen in terms of their positive minimum.  All these choices depend
only on \(\Delta\) and the pretarget slot system.

Choose
\[
 0<\zeta<\min\{1/48,m_{\rm high}/16\}.
\]
The frequency extension then costs at most
\(2\zeta<m_{\rm high}/8\).  Now choose the absolute line \(z_\infty>2\)
in Section~\ref{sec:tails} sufficiently far right that its large-row
sum has a saving larger than \(m_{\rm high}\) relative to \(C(\beta_*)\).
Its exponent depends only on
these real choices, so \(z_\infty\) also precedes the target.
Choose the pretarget cutoff needed for
Equation~\eqref{eq:shared-principal-nonvanishing}.  Any additional fixed
excluded primes required by Proposition~\ref{prop:probe-errors} after
the target is fixed may be added to \(S\): the same positive product-tail
majorant preserves the shared contraction, and neither \(T\) nor the
slot windows change.

For clarity, the nonprincipal real comparison is uniform over all
pointwise pairs in Lemma~\ref{lem:high-bin}.  Let
\(\epsilon_{\rm real}(d)\ge0\) collect the allocated real, mesh, and
moment losses for the relevant range.  The endpoint inequality,
the intermediate inequality, and the positive slopes give
\begin{equation}\label{eq:saving-criterion}
 E(d)+\epsilon_{\rm real}(d)
 \le(\beta_*-7/8)-\varepsilon_{\rm hi},
 \qquad \varepsilon_{\rm hi}>0,
\end{equation}
for every occurring \((R,q)\) and every moderate \(d\), with a fixed
fraction of \(m_{\rm high}\) still reserved.  Therefore the same
inequality holds for \(E_{\max,Z}(d)\), uniformly in \(Z\).
The small and large bounds have their stated separate savings.
Apply the subpower normalizer estimate once to the nonprincipal and
outer-row contributions, and absorb the remaining finite dyadic
multiplicities using a power below \(m_{\rm high}/8\).
Lemma~\ref{lem:principal-residue-interface} already includes its own normalizer allowance
\(\epsilon_{\rm pr}\).  The four high real allocations above cost less
than \(m_{\rm high}/2\); this last allowance leaves more than
\(3m_{\rm high}/8\).  The principal inequalities leave more than
\(3/4\) of each of \(m_w,m_z,m_P\).  Thus one may fix, before the target,
\[
 m=\frac14\min\{m_{\rm high},m_w,m_z,m_P\}>0
\]
as a common remaining high saving before retained height factors.
For the low estimate, apply Equation~\eqref{eq:part-II-normalized-low}
with the pretarget loss \(\Delta/2\), and set \(\omega=\Delta/2\).
Let \(\epsilon_{\rm ht}>0\) be the minimum of the finitely many detector
height allowances already chosen with these real losses.

Now fix a primitive target \(\eta\).  Choose its admissible fixed
arithmetic data and the final excluded set \(S\), including its conductor,
the common cutoff, and the further fixed exclusions just described.
The physical expression, \(H_\eta\), \(c_S\), and the signal all use this
same data.  The group \(T\) and the slot system remain independent of
the target.  Fix all internal moment, seminorm, and Sobolev orders
supplied by the preceding results for this datum.  Their uniformity over
moving moduli, common masks, and frozen labels gives a finite
\(A_\eta\) dominating the product of all retained high factors, including
the dyadic, prime, numerator-reflection, and row-count height factors.
It is independent of any later external test order.  The all-height
correction bounds and the direct global or absolute bounds in the shared
analytic lemmas likewise give a finite scale degree \(B_\eta\) for all
discarded physical pieces, independent of the later order \(N\).
Normalizing those pieces and summing their finitely many types only
enlarges this fixed \(B_\eta\).

We use the height ceiling verified at the end of
Section~\ref{sec:selection}.  Let \(A_{{\rm ht},\eta}\le A_\eta\)
dominate the finitely many detector height orders, enlarging \(A_\eta\)
if necessary, and put
\[
 \tau_{0,\eta}:=
       \frac{d_{\min}\epsilon_{\rm ht}}{20(1+A_{{\rm ht},\eta})}>0.
\]
This number is fixed after the retained orders for the target, but before
\(N\) and \(Z\).  If
\(0<\tau\le\min\{d_{\min}/100,\tau_{0,\eta}\}\) and \(T_1=Z^\tau\),
then, for \(U\ge Z^{d_{\min}}\) and sufficiently large \(Z\),
\[
 T_1\le U^{1/100},\qquad
 (1+T_1)^{A_{{\rm ht},\eta}}
 \le2^{A_{{\rm ht},\eta}}Z^{d_{\min}\epsilon_{\rm ht}/20}
 \le U^{\epsilon_{\rm ht}/10}.
\]
By the definition of \(\epsilon_{\rm ht}\), this proves every literal
condition \((1+T_1)^{A_{\mathcal A}}\le U^{\epsilon/10}\) in
Lemma~\ref{lem:detector-dyads} throughout the moderate range.  The first
bound also
makes the detector's crude error
\(U^{-189/100+o(1)}T_1^2\le U^{-187/100+o(1)}\) tend to zero.
Every buffered moderate range has \(d\le h+\zeta<1\); small and
absolute large rows use no buffered cutoff.

The fixed finite list of external coordinates, including the physical
\(\Im z\) coordinate, uses the single allocation in
Equation~\eqref{eq:stage-height-allocation}.  Its matrix is fixed by the
slot system and the finite retained estimates, and is independent of the
external derivative order.  For each such \(\tau\), the dyadic, witness,
and prime estimates may choose their auxiliary external orders after
\(\tau\), as their statements permit, so that their complete pointwise
bounds hold.  Choose a strict power gap when absorbing their discarded
parts; the resulting external seminorm constants are then absorbed by the
\(\tau\)-dependent lower threshold.  For any requested remaining order
\(N\), these auxiliary
orders may be increased further; this changes only external seminorm
constants and the lower threshold, not \(A_\eta\), \(B_\eta\), or
\(\tau_{0,\eta}\).  No internal moment is reapplied to an externally
differentiated profile.

Combining Lemma~\ref{lem:high-bin}, all the row ranges, and the principal
remainder estimate therefore gives, for every fixed \(N\) and sufficiently
large \(Z\) in this ceiling range,
\[
 |J_{\mathrm{II},\eta}(Z)-f_{\mathrm{II},\eta}(Z)|
 \ll_{\eta,N}
 Z^{C(\beta_*)-m}(1+T_1)^{A_\eta}
       +Z^{B_\eta}T_1^{-N}.
\]
The lower threshold is allowed to depend on \(\tau,N,\eta\).
This is Equation~\eqref{eq:late-height-hypothesis}.
Its low hypothesis is Equation~\eqref{eq:part-II-normalized-low}
with \(\omega=\Delta/2\), and the two functions do not contain \(T_1\).
Apply Lemma~\ref{lem:late-height-closure} with
\(\sigma_0=7/8\), \(C=C_{\mathrm{II}}\), this \(m\), and the verified
ceiling \(\tau_{0,\eta}\).  It chooses \(\tau_\eta\), then a common
dominating external order \(N_\eta\), and finally the threshold.
It gives the asserted high estimate with \(\sigma=m/2>0\).
Both \(\omega\) and \(\sigma\) were fixed before the target.
\end{proof}

Proposition~\ref{prop:parameter-order} and
Equation~\eqref{eq:shared-principal-nonvanishing} verify the hypotheses of
Proposition~\ref{lem:continuation-criterion} with
\(\sigma_0=7/8\), \(J_\eta=J_{\mathrm{II},\eta}\), and
\(f_\eta=f_{\mathrm{II},\eta}\).  This contradicts \(\beta_*>7/8\);
hence \(\beta_*\le7/8\).  Proposition~\ref{prop:hecke-dirichlet-transfer}
at \(\sigma_0=7/8\) extends the primitive Hecke conclusion to all
finite-order Hecke and Dirichlet \(L\)-functions, with the principal pole
allowed, and proves Theorem~\ref{thm:main}.
\begin{thebibliography}{99}

\bibitem{BaierBansal}
S. Baier and A. Bansal,
\emph{The large sieve with power moduli for $\mathbb Z[i]$},
International Journal of Number Theory \textbf{14} (2018), no. 10,
2737--2756.
\href{https://doi.org/10.1142/S1793042118501658}{doi:10.1142/S1793042118501658};
\href{https://arxiv.org/abs/1802.08964}{arXiv:1802.08964}.

\bibitem{BIMS}
V. Bhargava, G. Ivanyos, R. Mittal, and N. Saxena,
\emph{Irreducibility and deterministic $r$-th root finding over finite fields},
Proceedings of the 2017 ACM International Symposium on Symbolic and
Algebraic Computation (ISSAC 2017), ACM, 2017, 37--44.
\href{https://doi.org/10.1145/3087604.3087620}{doi:10.1145/3087604.3087620};
\href{https://www.cse.iitk.ac.in/users/nitin/papers/irred-and-rth-root.pdf}{author manuscript}.

\bibitem{BCSS}
K. Billington, M. Cheng, J. Schettler, and A. I. Suriajaya,
\emph{The average number of Goldbach representations and zero-free regions
of the Riemann zeta-function},
International Journal of Number Theory \textbf{21} (2025), no. 2, 289--316.
\href{https://doi.org/10.1142/S1793042125500150}{doi:10.1142/S1793042125500150};
\href{https://arxiv.org/abs/2306.09102}{arXiv:2306.09102}.

\bibitem{BGL}
V. Blomer, L. Goldmakher, and B. Louvel,
\emph{$L$-functions with $n$-th-order twists},
International Mathematics Research Notices \textbf{2014} (2014), no. 7,
1925--1955.
\href{https://doi.org/10.1093/imrn/rns257}{doi:10.1093/imrn/rns257};
\href{https://arxiv.org/abs/1112.1650}{arXiv:1112.1650}.

\bibitem{Clay}
Clay Mathematics Institute,
\emph{Riemann Hypothesis}, Millennium Prize Problems,
\url{https://www.claymath.org/millennium/riemann-hypothesis/}.

\bibitem{DFDS}
C. David, A. de Faveri, A. Dunn, and J. Stucky,
\emph{Non-vanishing for cubic Hecke $L$-functions},
preprint, arXiv:2410.03048v2 (2026).
\href{https://arxiv.org/abs/2410.03048}{arXiv:2410.03048}.

\bibitem{Dirichlet}
P. G. L. Dirichlet,
\emph{Beweis des Satzes, da\ss{} jede unbegrenzte arithmetische
Progression, deren erstes Glied und Differenz ganze Zahlen ohne
gemeinschaftlichen Factor sind, unendlich viele Primzahlen enth\"alt},
Abhandlungen der K\"oniglichen Preussischen Akademie der Wissenschaften
zu Berlin (1837), 45--81.
English translation by R. Stephan,
\href{https://arxiv.org/abs/0808.1408}{arXiv:0808.1408}.

\bibitem{DR}
A. Dunn and M. Radziwi\l\l,
\emph{Bias in cubic Gauss sums: Patterson's conjecture},
Annals of Mathematics \textbf{200} (2024), no. 3, 967--1057.
\href{https://doi.org/10.4007/annals.2024.200.3.3}{doi:10.4007/annals.2024.200.3.3};
\href{https://arxiv.org/abs/2109.07463}{arXiv:2109.07463}.

\bibitem{FG}
J. B. Friedlander and D. A. Goldston,
\emph{Sums of three or more primes},
Transactions of the American Mathematical Society \textbf{349} (1997),
no. 1, 287--310.
\href{https://doi.org/10.1090/S0002-9947-97-01652-8}{doi:10.1090/S0002-9947-97-01652-8}.

\bibitem{Galbraith}
S. D. Galbraith,
\emph{Mathematics of Public Key Cryptography},
Cambridge University Press, 2012.
\href{https://doi.org/10.1017/CBO9781139012843}{doi:10.1017/CBO9781139012843}.
Section and algorithm references use the author's extended and corrected
version 2.0, \href{https://www.math.auckland.ac.nz/~sgal018/crypto-book/ch2.pdf}{Chapter 2}.

\bibitem{GZ}
P. Gao and L. Zhao,
\emph{Moments of quadratic Hecke $L$-functions of imaginary quadratic
number fields},
Journal of Number Theory \textbf{209} (2020), 359--377.
\href{https://doi.org/10.1016/j.jnt.2019.09.002}{doi:10.1016/j.jnt.2019.09.002};
\href{https://arxiv.org/abs/1707.00091}{arXiv:1707.00091}.

\bibitem{GZSextic}
P. Gao and L. Zhao,
\emph{Moments and one level density of sextic Hecke $L$-functions},
Functiones et Approximatio Commentarii Mathematici \textbf{70} (2024),
no. 1, 7--28.
\href{https://doi.org/10.7169/facm/2057}{doi:10.7169/facm/2057};
\href{https://arxiv.org/abs/2201.01885}{arXiv:2201.01885}.

\bibitem{GL}
L. Goldmakher and B. Louvel,
\emph{A quadratic large sieve inequality over number fields},
Mathematical Proceedings of the Cambridge Philosophical Society
\textbf{154} (2013), no. 2, 193--212.
\href{https://doi.org/10.1017/S0305004112000370}{doi:10.1017/S0305004112000370};
\href{https://arxiv.org/abs/1112.1642}{arXiv:1112.1642}.

\bibitem{GM}
L. Guth and J. Maynard,
\emph{New large value estimates for Dirichlet polynomials},
Annals of Mathematics \textbf{203} (2026), no. 2, 623--675.
\href{https://doi.org/10.4007/annals.2026.203.2.6}{doi:10.4007/annals.2026.203.2.6};
\href{https://arxiv.org/abs/2405.20552}{arXiv:2405.20552}.

\bibitem{Hadamard}
J. Hadamard,
\emph{Sur la distribution des z\'eros de la fonction $\zeta(s)$ et ses
cons\'equences arithm\'etiques},
Bulletin de la Soci\'et\'e Math\'ematique de France \textbf{24} (1896),
199--220.
\href{https://doi.org/10.24033/bsmf.545}{doi:10.24033/bsmf.545}.

\bibitem{HBQuadratic}
D. R. Heath-Brown,
\emph{A mean value estimate for real character sums},
Acta Arithmetica \textbf{72} (1995), no. 3, 235--275.
\href{https://doi.org/10.4064/aa-72-3-235-275}{doi:10.4064/aa-72-3-235-275}.

\bibitem{HB}
D. R. Heath-Brown,
\emph{Kummer's conjecture for cubic Gauss sums},
Israel Journal of Mathematics \textbf{120}, part A (2000), 97--124.
\href{https://doi.org/10.1007/s11856-000-1273-y}{doi:10.1007/s11856-000-1273-y}.

\bibitem{Hecke}
E. Hecke,
\emph{Eine neue Art von Zetafunktionen und ihre Beziehungen zur
Verteilung der Primzahlen. Erste Mitteilung},
Mathematische Zeitschrift \textbf{1} (1918), 357--376;
\emph{Zweite Mitteilung}, Mathematische Zeitschrift \textbf{6} (1920),
11--51.
\href{https://doi.org/10.1007/BF01465095}{doi:10.1007/BF01465095}
(first paper).

\bibitem{Huxley}
M. N. Huxley,
\emph{The large sieve inequality for algebraic number fields},
Mathematika \textbf{15} (1968), no. 2, 178--187.
\href{https://doi.org/10.1112/S0025579300002540}{doi:10.1112/S0025579300002540}.

\bibitem{Kubota}
T. Kubota,
\emph{On automorphic functions and the reciprocity law in a number field},
Lectures in Mathematics, Department of Mathematics, Kyoto University,
no. 2, Kinokuniya Book-Store, Tokyo, 1969.
\url{http://hdl.handle.net/2433/84907}.

\bibitem{MP}
J. Maynard and K. Pratt,
\emph{Half-isolated zeros and zero-density estimates},
International Mathematics Research Notices \textbf{2024} (2024), no. 19,
12978--13014.
\href{https://doi.org/10.1093/imrn/rnae191}{doi:10.1093/imrn/rnae191};
\href{https://arxiv.org/abs/2206.11729}{arXiv:2206.11729}.

\bibitem{MilneCFT}
J. S. Milne,
\emph{Class Field Theory}, course notes, version 4.03, 6 August 2020.
\url{https://www.jmilne.org/math/CourseNotes/CFT.pdf}.

\bibitem{OpenAICubicFirstMoment}
OpenAI,
\emph{An unconditional first moment for cubic Gauss sums},
OpenAI Math Release preprint
\href{https://github.com/openai/math/blob/main/preprints/An-unconditional-first-moment-for-cubic-Gauss-sums-September-25-2026/paper.pdf}{OAI:An-unconditional-first-moment-for-cubic-Gauss-sums-September-25-2026},
2026.

\bibitem{Patterson}
S. J. Patterson,
\emph{A cubic analogue of the theta series},
Journal f\"ur die reine und angewandte Mathematik \textbf{296} (1977),
125--161.
\href{https://doi.org/10.1515/crll.1977.296.125}{doi:10.1515/crll.1977.296.125}.

\bibitem{Riemann}
B. Riemann,
\emph{Ueber die Anzahl der Primzahlen unter einer gegebenen Gr\"osse},
Monatsberichte der K\"oniglichen Preussischen Akademie der Wissenschaften
zu Berlin (1859), 671--680.
Manuscript, transcription, and translation at
\url{https://www.claymath.org/library/historical/riemann/r.html}.

\bibitem{TaoNonresidues}
T. Tao,
\emph{The Elliott--Halberstam conjecture implies the Vinogradov least
quadratic nonresidue conjecture},
Algebra \& Number Theory \textbf{9} (2015), no. 4, 1005--1034.
\href{https://doi.org/10.2140/ant.2015.9.1005}{doi:10.2140/ant.2015.9.1005};
\href{https://msp.org/ant/2015/9-4/ant-v9-n4-p09-p.pdf}{publisher PDF}.

\bibitem{TZ}
J. Thorner and A. Zaman,
\emph{A unified and improved Chebotarev density theorem},
Algebra \& Number Theory \textbf{13} (2019), no. 5, 1039--1068.
\href{https://doi.org/10.2140/ant.2019.13.1039}{doi:10.2140/ant.2019.13.1039};
\href{https://arxiv.org/abs/1803.02823}{arXiv:1803.02823}.

\bibitem{ValleePoussin}
C.-J. de la Vall\'ee Poussin,
\emph{Recherches analytiques sur la th\'eorie des nombres premiers.
Premi\`ere partie: La fonction $\zeta(s)$ de Riemann et les nombres
premiers en g\'en\'eral},
Annales de la Soci\'et\'e scientifique de Bruxelles \textbf{20},
deuxi\`eme partie (1896), 183--256.
Reprinted in \emph{Recherches analytiques sur la th\'eorie des nombres
premiers}, Hayez, Brussels, 1897,
\url{https://books.google.com/books?id=7e0GAAAAYAAJ}.

\end{thebibliography}
\end{document}
